% =========================================
% CHAPTER 7: SECURITY ANALYSIS
% =========================================
\chapter{Security Analysis and Failure Modes}
\label{ch:security}

This chapter examines the security properties of both rollup architectures under adversarial conditions. We define a threat model, analyze the strengths and vulnerabilities of each architecture, and evaluate the impact of L1 chain reorganizations on L2 finality.

\section{Threat Model}

We consider five adversarial scenarios that are relevant to any Layer~2 deployment:

\begin{enumerate}
    \item \textbf{Malicious sequencer}: the sequencer attempts to commit invalid state transitions to L1.
    \item \textbf{L1 chain reorganization}: one or more confirmed L1 blocks are replaced, invalidating L2 commitments anchored to those blocks.
    \item \textbf{Sequencer failure}: the sequencer crashes or becomes temporarily unavailable.
    \item \textbf{Censorship attack}: the sequencer selectively excludes certain transactions from batches.
    \item \textbf{Data unavailability}: L2 state data is not published or becomes inaccessible, preventing independent verification.
\end{enumerate}

The following sections analyze how each architecture responds to these threats.

\section{ZK Rollup Security}

\subsection{Strength: Cryptographic Validity}

The security of a ZK~Rollup derives from the mathematical properties of the underlying proof system:
\begin{itemize}
    \item \textbf{Soundness}: a valid proof guarantees that the corresponding state transition is correct; it is computationally infeasible to forge a proof for an invalid transition.
    \item \textbf{Zero-knowledge}: the proof reveals no information about the individual transactions within the batch, only that the transition is valid.
    \item \textbf{Non-interactivity}: proof verification is a single-step operation that requires no further communication between the prover and the verifier.
\end{itemize}

The practical implication is that any invalid state is rejected at the point of L1 verification. Invalid commitments never enter L1 storage, providing a strong safety guarantee.

\subsection{Vulnerability: Prover Availability}

If the sequencer produces valid transactions but fails to generate proofs (whether due to a crash, resource exhaustion, or deliberate refusal), the system enters a degraded state:
\begin{itemize}
    \item transactions continue to be accepted at the L2 level but are not finalized on L1;
    \item users cannot prove the current L2 state to external parties or to L1;
    \item the system is effectively \textit{censored but not compromised}---no funds are at risk, but progress halts.
\end{itemize}

The standard mitigation is a force-inclusion mechanism that requires the sequencer to produce proofs within a deadline or face penalties (such as collateral slashing or sequencer replacement).

\subsection{Vulnerability: Proof System Compromise}

If the zero-knowledge proof system itself contains a mathematical flaw, it may become possible to construct proofs for invalid state transitions. This is a low-probability but high-impact risk. Standard mitigations include:
\begin{itemize}
    \item adopting well-studied, battle-tested proof systems such as PLONK or STARK;
    \item commissioning regular security audits by cryptography experts;
    \item monitoring the academic literature for newly discovered vulnerabilities in the underlying assumptions.
\end{itemize}

\section{Optimistic Rollup Security}

\subsection{Strength: Economic Incentives}

The security of an Optimistic~Rollup is grounded in economic game theory rather than cryptographic proofs:
\begin{itemize}
    \item \textbf{Collateral requirements}: both the sequencer and any challenger must post bonds before participating in the protocol.
    \item \textbf{Slashing}: a party found to be dishonest (either a malicious sequencer or a frivolous challenger) loses its posted collateral.
    \item \textbf{Rational incentives}: under standard game-theoretic assumptions, honest behavior is the economically dominant strategy.
\end{itemize}

This design allows invalid states to be challenged and reverted, but only through an economic dispute process rather than an automatic cryptographic check.

\subsection{Vulnerability: Challenger Absence}

The most significant vulnerability of an Optimistic~Rollup is the possibility that no honest challenger submits a fraud proof during the challenge period. This can occur under the following conditions:
\begin{itemize}
    \item low network participation, with few nodes running full archival copies of the L2 state;
    \item prohibitively high challenge costs that exceed the economic incentive to challenge;
    \item coordinated collusion between the sequencer and a majority of potential challengers.
\end{itemize}

We can model the probability of a successful attack (i.e., no challenge) as follows. Let~$p$ denote the probability that any single observer detects the fraud, and let~$n$ denote the number of independent observers. Then:
\begin{equation}
P(\text{no challenge}) = (1 - p)^n
\end{equation}

For example, with $p = 0.1$ and $n = 10$:
\begin{equation}
P(\text{no challenge}) = (0.9)^{10} \approx 0.35
\end{equation}

This illustrative calculation shows that even with 10 independent observers, each having a 10\% chance of detecting fraud, there remains a 35\% probability that no challenge is submitted. In practice, production systems mitigate this through dedicated watchtower services and by increasing the number of independent validators.

\subsection{Vulnerability: Withdrawal Latency}

Users who wish to withdraw funds from an Optimistic~Rollup to L1 must wait for the full challenge period (7~days in production) to elapse. This has several practical consequences:
\begin{itemize}
    \item users who need faster access to their funds must rely on third-party liquidity providers, who charge a fee for the service;
    \item withdrawals of small amounts become disproportionately expensive relative to liquidity provider fees;
    \item the user experience is significantly worse than that of a ZK~Rollup, where withdrawals can be processed as soon as the proof is verified on L1.
\end{itemize}

Proposed mitigations include liquidity networks and hybrid models that use ZK proofs to accelerate the withdrawal process.

\section{L1 Reorganization Impact}

\subsection{Scenario: Shallow Reorg (1~Block)}

Consider the following situation: an L1 block at height~$n$ contains a ZK proof~$\pi_1$ for batch~$B_1$. A 1-block reorganization occurs, and block~$n$ is replaced by a new block that does not contain~$\pi_1$.

The impact on the ZK~Rollup is as follows:
\begin{itemize}
    \item the proof $\pi_1$ is no longer part of the L1 canonical chain;
    \item the L2 state update that was based on $\pi_1$ becomes invalid;
    \item the system reverts to the previous valid state;
    \item the sequencer must resubmit the batch to L1.
\end{itemize}

The recovery process is automatic:
\begin{enumerate}
    \item The sequencer detects the reorganization by monitoring the L1 chain.
    \item It regenerates the proof (possibly with updated batch ordering if necessary).
    \item It resubmits the proof to L1.
    \item The system re-finalizes once the new proof is included and confirmed.
\end{enumerate}

The expected recovery time can be estimated as:
\begin{equation}
T_{\text{recovery}} = T_{\text{detection}} + T_{\text{reproof}} + T_{\text{re-inclusion}}
\end{equation}

For a 1-block reorg on Ethereum: detection~$\approx$~12s, reproof~$\approx$~95ms, re-inclusion~$\approx$~12s, yielding a total recovery time of approximately \textbf{24~seconds}. Importantly, no data is lost during this process.

\subsection{Scenario: Deep Reorg ($>$7~Blocks)}

A deep reorganization ($>$7~blocks) is an extremely unlikely event on Ethereum, particularly following the transition to proof-of-stake and the introduction of finality gadgets. Such an event would require either a successful 51\% attack or a catastrophic network partition.

If a deep reorg were to occur:
\begin{itemize}
    \item both ZK and Optimistic~Rollups would need to resubmit multiple batches or state roots;
    \item in the Optimistic case, economic incentives could break down if challengers are unable to participate during the disruption;
    \item the system would likely require manual intervention or social consensus to recover.
\end{itemize}

For practical purposes, deep reorgs are not considered a realistic threat for established L1 networks.

\section{Comparative Security Summary}

Table~\ref{tab:security-summary} provides a qualitative summary of the security properties of each architecture across the dimensions considered in this chapter.

\begin{table}[h]
    \centering
    \begin{tabular}{l|c|c}
        \hline
        \textbf{Security Property} & \textbf{ZK Rollup} & \textbf{Optimistic} \\
        \hline
        Cryptographic safety & Strong & Economic \\
        \hline
        Recovery speed & Moderate & Moderate \\
        \hline
        Censorship resistance & Good & Medium \\
        \hline
        Liveness requirement & Prover only & Prover + challenger \\
        \hline
        Finality speed & Fast ($\sim$15s) & Slow (7~days) \\
        \hline
        User exit cost & Low & High (7~days or fee) \\
        \hline
    \end{tabular}
    \caption{Comparative security summary}
    \label{tab:security-summary}
\end{table}
